% Program synthesis paper — map of the Infinity Compression paper suite and Lean artifact.
\documentclass[11pt]{article}
\input{../suite_preamble}
\input{../suite_macros}

\title{Certification, Realization, and Obstruction:\\[0.35em]
\large A Universal Fiber Architecture}
\author{Nova Spivack}
\date{2026}

\begin{document}

\maketitle

\begin{abstract}
This paper closes the editorial loop for a multi-manuscript program in machine-checked
mathematics. The executive claim is that a single \emph{collapse / refinement / fiber}
architecture---certification at a bare carrier, strict comparison upstairs, obstruction
organized as fiber and section data---is instantiated internally in dependent type theory,
validated on twelve external mathematical families, and then discharged into algebra,
arithmetic (embedding problems), topology (Quillen for Galois connections), and
logic/computability (self-certification and halting anchors). We state the role of each
companion paper, a cross-domain dictionary, route labels used in the program (C/A/B/D),
and what remains open. The intended reader is an advanced mathematician who has not read
the series; citations point to the formal manuscripts and the Lean repository.
\end{abstract}

\keywords{formal verification, certification vs.\ realization, fibers, obstruction,
  Infinity Compression, Lean~4, program architecture}

\tableofcontents
\newpage

\section{Executive statement}

\textbf{One-sentence thesis.} Canonical certification at a collapsed bare carrier does not
exhaust structurally significant information: the gap is organized by a forgetful map,
nontrivial fibers, and explicit obstructions---and that pattern repeats across internal
formalization, external benchmarks, and classical mathematical domains.

\section{Certification, realization, and non-exhaustion}

\paragraph{Certification.}
A \emph{certification} layer is a theorem-level invariant packaged as a canonical record
(proof-irrelevant at that carrier): many routes may exist, but the official statement
collapses to a unique bare certificate.

\paragraph{Realization.}
\emph{Realization} refers to proof-relevant or structured data above that collapse:
witnesses, roles, sections, or fiber points that are mathematically load-bearing once the
comparison map is fixed.

\paragraph{Non-exhaustion.}
\emph{Non-exhaustion} is the proved claim that certification does not determine realization:
distinct legitimate data can share the same bare image, organized by $\pi$ and fibers.

\section{Flagship internal theorem}

The Infinity Compression flagship paper~\cite{SpivackIC} is the \emph{internal origin
theorem}: collapse of standard extraction, reflective split, strict refinement of the
typed forgetful map, and partial fiber classification---all in the library's native
carriers.

\section{External validation (transferability)}

The twelve-tranche survey~\cite{SpivackValidation} establishes \emph{transferability}:
the same schema applies to quotients, products, dependent sums, coproducts, subtypes,
orbit and ideal quotients, pushouts, classifiers, localization, pullbacks, and a
non-surjective control. Its job is not theorem-level leverage in classical fields; it
rules out idiosyncrasy.

\section{Algebraic and arithmetic engine}

The group-extension paper~\cite{SpivackFiber} is the \emph{algebraic and arithmetic engine}:
section cocycles, splitting iff trivial cocycle, embedding problems as extension splitting,
local-to-global obstruction, and the documented bridge toward Galois-theoretic towers.

\section{Topological discharge}

Quillen's Theorem~A for Galois connections~\cite{SpivackQuillen} is the \emph{topology
discharge}: nerves, contractible fibers over down-sets, and homotopy-equivalence data at
the simplicial level for the Galois-connection case.

\section{Logic and computability anchor}

Route~D in the program supplies a computability-facing discharge: self-certification
phenomena aligned with standard undecidability infrastructure (e.g.\ halting and Rice
scaffolding in the Mathlib \texttt{ComputablePred} layer), making the analogy to
classical incompleteness precise without identifying the main theorem with G\"odel's
syntax. Exact statements live in the Lean \texttt{RouteD/} and \texttt{Summit/} modules;
this synthesis paper does not restate every definition.

\section{Cross-domain dictionary}

The following table is a navigational index, not a formal commuting diagram.

\begin{center}
\small
\begin{tabular}{@{}p{0.22\textwidth}p{0.34\textwidth}p{0.36\textwidth}@{}}
\toprule
\textbf{Domain} & \textbf{Carrier intuition} & \textbf{Obstruction / fiber role} \\
\midrule
Internal IC & Bare summit certificate vs.\ enriched split & Non-injectivity of $\pi_A$; canonical fiber \\
External T1--T12 & Baseline vs.\ structured bundle & Fibers, sections, right inverses \\
Algebra & Group extension & Section cocycle; splitting \\
Arithmetic bridge & Embedding problem & Solvable iff trivial cocycle \\
Topology & Galois connection $l \dashv u$ & Contractible down-set fibers \\
Computability & Self-certification predicates & Undecidability anchors \\
\bottomrule
\end{tabular}
\end{center}

\section{Routes C, A, B, and D (program labels)}

\begin{itemize}[leftmargin=2em]
  \item \textbf{Route C (headline):} reflective non-exhaustion at the flagship carriers.
  \item \textbf{Route A (infrastructure):} supporting general lemmas (e.g.\ collision
    patterns, forgetful infrastructure) used by the headline.
  \item \textbf{Route B (applications):} algebra, topology, arithmetic bridges---the
    general-method papers.
  \item \textbf{Route D (logic/computability):} halting and Rice-style scaffolding aligned
    with computability theory rather than informal metaphor.
\end{itemize}

\section{Why this is a reusable research instrument}

The architecture is portable because it isolates \emph{where} collapse occurs, \emph{where}
comparison must move upstairs, and \emph{what} invariant obstructs lifting---independent
of any single application domain. Lean provides a checkable artifact; the paper suite
separates internal origin, transferability, domain discharges, and summit-level claims.

\section{What remains open}

Full fiber classification in the flagship setting; complete simplicial homotopy
formalization behind Quillen; full $H^1$/$H^2$ classification in Mathlib; stronger
unified quantification across all carriers in a single declaration; publication and Mathlib
submission scheduling.

\section*{Acknowledgments}

Repository: \url{https://github.com/novaspivack/infinity-compression-lean}.

\bibliographystyle{alpha}
\begin{thebibliography}{99}

\bibitem{SpivackIC}
N.~Spivack.
\newblock \emph{Canonical Certification Does Not Exhaust Reflective Structure:
  Reflective Split, Strict Refinement, and Fiber Structure in Infinity Compression}.
\newblock Companion manuscript, 2026.

\bibitem{SpivackValidation}
N.~Spivack.
\newblock \emph{External Validation of a Positive-Closure Proof Architecture:
  A Multi-Tranche Survey}.
\newblock Companion manuscript, 2026.

\bibitem{SpivackFiber}
N.~Spivack.
\newblock \emph{Fiber Architecture for Group Extensions in Lean~4: Cocycles, Splitting,
  and the Cohomological Bridge}.
\newblock Companion manuscript, 2026.

\bibitem{SpivackQuillen}
N.~Spivack.
\newblock \emph{Quillen's Theorem~A for Galois Connections: A Machine-Checked
  Formalization in Lean~4}.
\newblock Companion manuscript, 2026.

\bibitem{SpivackSummit}
N.~Spivack.
\newblock \emph{Reflective Non-Exhaustion: Certification, Realization, and the Summit Program}.
\newblock Program manuscript, 2026.

\end{thebibliography}

\end{document}
